\documentclass{article}

\usepackage[utf8]{inputenc}     % pdfLaTeX
\usepackage[T1]{fontenc}
\usepackage[magyar]{babel}

\title{A Középkor Irodalma}
\author{Lengyel Zoltán Péter, Rovenszky Robin}
\date{Legutóbb frissítve: 2026. Április 13.}

\begin{document}

\maketitle
\tableofcontents
\newpage

\section{A középkor meghatározása}

\begin{itemize}
    \item \emph{media aetas}: reneszánsz ember nevezte el, köztes kor az antikvitás és a reneszász között
    \item Történelmi középkor 476-1492 (Nyugat Római birodalom bukása - Kolumbusz Kristóf felfedezi Amerikát)
    \item Művészettörténeti középkor: 2.sz-13/14. sz
\end{itemize}

\section{A középkor kultúrája}
\begin{itemize}
    \item Létszemléleti, vallásfilozófiai alap: keresztény vallás
    \item Transzcendes gondolkodás
    \item Ciklikus világkép
    
    \item Univerzalizmus 
    \begin{itemize}
        \item Isten egyetemessége
        \item Katolikus egyház egyetemessége
        \item Latin nyelv egyetemessége
    \end{itemize}

    \item Személyiségkép: közösségben létezés (nincs individualizmus) aszkézis, anonimitás

    \item Kettős világkép 
    \begin{itemize}
        \item Evilág 
        \begin{itemize}
            \item Földi szféra
            \item Múlandó,  véges
            \item Megismerhető, megtapasztalható
            \item Túlvilág
        \end{itemize}
        
        \item Túlvilág 
        \begin{itemize}
            \item Transzcendens
            \item Megtapasztalhatatlan
            \item Örök, állandó
        \end{itemize}
    \end{itemize}

    \item Kettős emberkép \begin{itemize}
        \item Szellemi (lélek)
        \item Anyagi (test)
    \end{itemize}

\end{itemize}

\section{Őskeresztény irodalom}
\begin{itemize}
    \item Biblia
    \item Apróbb legendák
    \item Államvallás
    \item 313 - Milánói ediktum (szabad vallásgyakorlat a keresztényeknek)
    \item 325 - \textbf{níceai zsinat} - Szentháromság kimondása
    \item 391 - Államvallás lett a kereszténység
\end{itemize}

\section{Kora-középkori egyházi irodalom}
\begin{itemize}
    \item Latin nyelvű
    \item Prédikációk, templomi énekek, himnuszok, siratóénekek
    
    \item A középkort a vallás határozza meg. DE! \textbf{szerzetesrendek}, \emph{engedelmesség, tisztaság, szegénység}: 3-as fogadalom 
    \begin{itemize}
        \item Ferencesek - élén Szent Ferenc
    \end{itemize}
\end{itemize}

\subsection{Szent Ferenc}
\begin{itemize}
    \item Assisi Szent Ferenc
    \item Itáliában, Assisiben született
    \item Apja posztókereskedő
    \item Egy betegség hatására tért meg
    \item Elment Római zarándoklatra
    \item Ruhát cserélt egy koldussal
    \item Apja rajtakapja hogy koldul az utcán
    \item Apja feljelentette a saját fiát 
    \item Meztelenre vetkőzött a város közepén 
    \item Teljesen megtért, Istennek szentelte életét, 3 templomot állított helyre
    \item Kolduló rendet hoz létre $\rightarrow$ a szegénységi fogadalom a legfontosabb, nem kell sok az élethez
    \item Aszkéta (önmegtartóztató) életforma
    \item Rómában bemutatják III. Ince pápának, aki elismerte a rendet
    \item Megalakul a női ág: Szent Klára - klarisszák rendje
    \item Nem csak alamizsnát kértek, dolgoztak is (betegápolás, földművelés ... stb)
\end{itemize}

\textbf{Assisi Szent Ferenc: Naphimnusz}

\end{document}